\section{Integral families for the native cost comparisons}
\label{app:families}
The denominator convention is $D_i=q_i^2-m_i^2+\mathrm{i}0$, with the measure in Eq.~\eqref{eq:family}. This appendix collects the ordered momenta used by the later cost studies so that the notes specify their physical inputs independently of the companion paper. Appendix~\ref{app:adaptive} gives the masses, kinematic profiles, and target indices for the new comparisons. Appendix~\ref{app:higher-scattering} adds the connected higher-loop scattering families.

The sunrise uses
\begin{equation}
\boldsymbol q_{\rm S}=(\ell_1,\ell_2,p-\ell_1-\ell_2,p-\ell_1,p-\ell_2).
\end{equation}
Its first three entries are physical denominators and its last two entries are auxiliary. Their indices remain nonpositive.

\subsection{Native multiloop families}

The sunrise comparisons allow independent masses on the same three physical propagators. The other native comparison families use
\begin{align}
 \boldsymbol q_{\rm DB}={}&(\ell_1,\ell_1-p_1,\ell_1-p_1-p_2,
 \ell_2,\ell_2-p_1-p_2,\ell_2-p_1-p_2-p_3,\nonumber\\
 &\ell_1-\ell_2,\ell_1-p_3,\ell_2-p_1),\\
 \boldsymbol q_{\rm V}={}&(\ell_1,\ell_2,\ell_3,\ell_1-\ell_2,
 \ell_2-\ell_3,\ell_1-\ell_3),\\
 \boldsymbol q_{\rm L}={}&(\ell_1,\ell_1-p,\ell_2,\ell_2-p,\ell_3,
 \ell_3-p,\ell_1-\ell_2,\ell_2-\ell_3,\ell_1-\ell_3).
\end{align}
The double box has two loops and seven physical denominators; its last two entries are auxiliary. The vacuum has three loops and six physical denominators. The ladder has three loops and eight physical denominators; its last entry is auxiliary. Auxiliary indices remain nonpositive. The massless double box has $p_i^2=0$ and uses $2p_1\cdot p_2=s$, $2p_1\cdot p_3=t$, $2p_2\cdot p_3=-s-t$. This local definition of $t$ differs from the initial box convention because the routing differs.


\subsection{Scattering structures and symbolic invariants}

The massless nonplanar box uses the routing of Ref.~\cite{KiraOriginal}:
\begin{align}
 \boldsymbol q_{\rm NP}={}&(\ell_1,\ell_2,p_1-\ell_1,p_2-\ell_2,
 p_1+p_2-p_3-\ell_1,\nonumber\\
 &p_1-\ell_1+\ell_2,p_1-p_3-\ell_1+\ell_2,
 \ell_1-p_3,\ell_2-p_1-p_2).
\end{align}
The first seven entries are physical propagators, and the final two are auxiliary entries. Here $p_1^2=p_2^2=p_3^2=(p_1+p_2-p_3)^2=0$, $s=(p_1+p_2)^2$ and $t=(p_1-p_3)^2$. Expanding these squares gives $2p_1\cdot p_2=s$, $2p_1\cdot p_3=-t$ and $2p_2\cdot p_3=s+t$.

The massless pentagon uses $\sum_{i=1}^5p_i=0$, $p_i^2=0$ and
\begin{equation}
 \boldsymbol q_{\rm P}=(\ell,\ell-p_1,\ell-p_1-p_2,
 \ell-p_1-p_2-p_3,\ell-p_1-p_2-p_3-p_4).
\end{equation}
All five denominators are physical. The independent invariants are $(s_{12},s_{23},s_{34},s_{45},s_{51})$, where $s_{ij}=(p_i+p_j)^2$. Squaring $p_1+p_2+p_3=-(p_4+p_5)$ gives $2p_1\cdot p_3=s_{45}-s_{12}-s_{23}$; the other nonadjacent products follow by cyclic relabelling. The box with an external virtuality uses $\boldsymbol q=(\ell,\ell-p_1,\ell-p_1-p_2,\ell-p_3)$, $p_1^2=p_2^2=p_3^2=0$, and $v=(p_1+p_2-p_3)^2$. With the nonplanar convention for $s,t$, expanding $v$ gives $2p_2\cdot p_3=s+t-v$.
